\newpage
\section{Determinant of a Square Matrix}
Let $A = [a_{ij}]\in M_n(\mathbb{R})$. The determinant of $A$ is defined as follows:
\begin{itemize}
    \item If $n = 1$, then $\det A = a_{11}$
    \item Suppose $n \geq 1$ and assume that the determinant of a matrix of size $n - 1$ is defined. Let $M_{ij}$ be the matrix obtained from $A$ by deleting row $i$ and column $j$ of $A$. Define
    $$C_{ij} = (-1)^{i+j} \det M_{ij}$$ wherein $(i, j)$ is the cofactor of $A$.
\end{itemize}

Then we define the determinant as a cofactor expansion along row $i$,
$$\det A = a_{i1}C_{i1} + a_{i2}C_{i2} + \cdots + a_{in}C_{in}$$

Similarly, we could also define it as a cofactor expansion along column $j$,
$$\det A = a_{1j}C_{1j} + a_{2j}C_{2j} + \cdots + a_{nj}C_{nj}$$

\lline

\begin{example}
Obtain a formula for the determinant of $\displaystyle A = \left[ 
\begin{matrix}
    a_{11} & a_{12} \\
    a_{21} & a_{22}
\end{matrix}
\right]$
\end{example}

\textit{Solution.} By cofactor expansion along row 1,
\begin{align*}
\det A &= a_{11}C_{11} + a_{12}C_{12} \\
&=a_{11}(-1)^{1+1}\det [a_{22}] + a_{12}(-1)^{1+2}\det [a_{21}] \\
&=\boxed{a_{11}a_{22} - a_{12}a_{21}}
\end{align*}

\begin{example}
Find the determinant of $\displaystyle B = \left[
\begin{matrix}
    5 & 6 \\
    4 & 8
\end{matrix}
\right]$.
\end{example}

\textit{Solution.} Recall that from the previous example, we have obtained a determinant formula for a square matrix of size 2, we could use that to solve for $\det B$
$$\det B = 5(8) - 6(4) = \boxed{16}$$

\begin{example}
Find the determinant of $\displaystyle B = \left[ 
\begin{matrix}
    3 & 1 & -4 \\
    2 & 5 & 6 \\
    1 & 4 & 8
\end{matrix}
\right]$.
\end{example}

\textit{Solution.} By cofactor expansion along column 1,
$$\det B = 3(-1)^{1+1}
\vmat{
5 & 6 \\
4 & 8
} + 
2(-1)^{2 + 1}
\vmat{
1 & -4 \\
4 & 8
} +
1(-1)^{3 + 1}
\vmat{
1 & -4 \\
5 & 6
}
$$
$$=3(16) -2(24) + 26 = \boxed{26}$$

\newpage
\begin{example}
Find the determinant of $\displaystyle C = \bmat{
2 & 3 & 0 & 0 \\
1 & 1 & 0 & 0 \\
4 & 6 & 1 & -5 \\
0 & 5 & 3 & 0
}$.
\end{example}

\textit{Solution.} Notice that column 4 has 3 zeros which could help us compute for the determinant for efficiently,
$$\det C = 0 + 0 -5(-1)^{3 + 4}\vmat{
2 & 3 & 0 \\
1 & 1 & 0 \\
0 & 5 & 3
} + 0 = 5 \vmat{
2 & 3 & 0 \\
1 & 1 & 0 \\
0 & 5 & 3
}$$
We will perform cofactor expansion along column 3 since it has the least amount of zeros,
$$=5\left(3(-1)^{3+3}\vmat{
2 & 3 \\
1 & 1
}\right) = 15(2 - 3) = \boxed{-15}$$

\begin{example}
The examples that will be shown here will help us determine some properties of determinants. Let $\displaystyle A = \bmat{
a & b \\
c & d
}$, and $\displaystyle B = \bmat{
a_{11} & a_{12} & a_{13} \\
0 & a_{22} & a_{23} \\
0 & 0 & a_{33}
}$
\begin{itemize}
    \item Suppose we switch row 1 and row 2 of matrix $A$, then the determinant is
    $$\vmat{
    c & d \\
    a & b
    } = bc - ad = -(ad -bc) = -\det A$$
    \item Suppose we add a constant $r \in \mathbb{R}$ to row 2 of matrix $A$, then the determinant is
    $$\vmat{
    a & b\\
    rc & rd
    } = ard - brc = r(ad - bc) = r\det A$$
    \item Suppose we multiply $r$ to row 1 of matrix $A$ and then add it to row 2, then the determinant is
    $$\vmat{
    a & b\\
    ra + c & rb + d
    } = a(rb + d) - b(ra + c) = ad - bc = \det A$$
    \item Suppose we perform cofactor expansion along row 3 of matrix $B$ to get its determinant,
    $$\det B = a_{33}(-1)^{3 + 3}\vmat{
    a_{11} & a_{12} \\
    0 & a_{22}
    } = a_{33}(a_{11}a_{22}) = a_{11}a_{22}a_{33}$$
\end{itemize}
\end{example}

\lline

\underline{\textbf{Some Properties of Determinant}}. Let $A, B \in M_n(\mathbb{R})$,

1. If $B$ is obtained from $A$ by:
\begin{itemize}
    \item switching two rows/columns of $A$ then $\det B = -\det A$
    \item multiplying $r \in \mathbb{R}$ to a row/columns of $A$ then $\det B = r\det A$

    e.g. $\det(rA) = \underbrace{r \times r \times \cdots \times r}_{n}\det A = r^n\det A$
    \item adding a multiple of a row/column of $A$ to another row/column of $A$ then $\det B = \det A$
\end{itemize}

2. $\det(A^{\mathsf{T}}) = \det A$

3. If $A$ has a row/column full of 0s then $\det A = 0$.

4. If $A$ has two equal rows/columns, then $\det A = 0$.

5. If $A = [a_{ij}]$ is upper or lower triangular, then $\det A = a_{11}a_{22}\cdots a_{nn}$, which is the product of all the main diagonal entries of the matrix.

6. $\det(AB) = \det A \det B$, and in particular, $\det(A^k) = (\det A)^k$ for all $k \in \mathbb{Z}^+$.

\lline

\begin{theorem}
If $E \in M_n(\mathbb{R})$ is elementary, then $\det E \neq 0$.
\end{theorem}

\begin{theorem}
$A \in M_n(\mathbb{R})$ is nonsingular iff $\det A \neq 0$, in this case
$$\det(A^{-1}) = \frac{1}{\det A}$$
\end{theorem}

\begin{definition}
Let $A \in M_n(\mathbb{R})$ and $C_{ij}$ be the $(i, j)$ column of $A$. The \textbf{adjoint} of $A$ is $\adj A = [C_{ij}]^{\mathsf{T}}$.
\end{definition}

\begin{example}
Find the adjoint of the following:

1. $\displaystyle A = \bmat{a & b \\ c & d}$

\textit{Solution.}

\begin{align*}
C_{11} = (-1)^{1 + 1}d = d \quad & C_{21} = (-1)^{2 + 1} b = -b \\
C_{12} = (-1)^{1 + 2}c = -c \quad & C_{22} = (-1)^{2 + 2}a = a
\end{align*}

$$\therefore \adj A = \bmat{ d & -b \\ -c & a}$$

2. $\displaystyle B = \bmat{1 & 1 & 0 \\ 2 & 3 & 4 \\ 5 & 8 & 9}$
\end{example}

\textit{Solution.} \label{example:4.6}
\begin{align*}
C_{11} &= \vmat{3 & 4 \\ 8 & 9} = -5 & C_{21} &= -\vmat{1 & 0 \\ 8 & 9} = -9 & C_{31} &= \vmat{1 & 0 \\ 3 & 4} = 4 \\
C_{12} &= -\vmat{2 & 4 \\ 5 & 9} = 2 & C_{22} &= \vmat{1 & 0 \\ 5 & 9} = 9  & C_{32} &= -\vmat{1 & 0 \\ 2 & 4} = -4 \\
C_{13} &= \vmat{2 & 3 \\ 5 & 8} = 1  & C_{23} &= -\vmat{1 & 1 \\ 5 & 8} = -3 & C_{33} &= \vmat{1 & 1 \\ 2 & 3} = 1
\end{align*}
$$\therefore \adj B = \bmat{-5 & -9 & 4 \\ 2 & 9 & -4 \\ 1 & -3 & 1}$$

\begin{theorem}[Adjoint Formula]
If $A \in M_n(\mathbb{R})$ then,
$$A(\adj A) = (\det A)I_n = (\adj A)A$$
If $A$ is nonsingular then,
$$A\left(\frac{1}{\det A} \cdot \adj A\right) = I_n$$
$$A^{-1} = \frac{1}{\det A} \adj A$$
\end{theorem}

\begin{example}
If $\displaystyle B = \bmat{
a & b \\
c & d
}$ is nonsingular then,
$$A^{-1} = \frac{1}{ad - bc}\bmat{d & -b \\ -c & a}$$
\end{example}

\begin{example}
Find the inverse of $\displaystyle B = \bmat{1 & 1 & 0 \\ 2 & 3 & 4 \\ 5 & 8 & 9}$.
\end{example}

\textit{Solution.}

Recall that if we perform Type III operations in any row or column of a matrix, then its determinant does not change. Thus, we could multiply -1 to the the first column and add that to the second column, we would get the following matrix:
$$\bmat{1 & 0 & 0 \\ 2 &1 & 4 \\ 5 & 3 & 9} \Longrightarrow \det B = \vmat{1 & 0 & 0 \\ 2 & 1 & 4 \\ 5 & 3 & 9}$$

Let us perform cofactor expansion along row 1,
$$\det B = 1 \vmat{1 & 4 \\ 3 & 9} = 9 - 12 = -3$$

Since we have already obtained $\adj B$ from \hyperref[example:4.6]{Example 4.6}, we could use that to find the inverse using the adjoint formula,
$$B^{-1} = -\frac{1}{3}\bmat{-5 & -9 & 4 \\ 2 & 9 & -4 \\ 1 & -3 & 1}=\boxed{\bmat{5/3 & 3 & -4/3 \\ -2/3 & -3 & 4/3 \\ -1/3 & 1 & -1/3}}$$

\begin{theorem}[Cramer's Rule]
Let $A \in M_n(\mathbb{R})$ be nonsingular and $b \in M_{n \times 1}(\mathbb{R})$. The matrix equation $Ax = b$ has a unique solution $\displaystyle x = \bmat{x_1 & x_2 & \cdots & x_n}^{\mathsf{T}}$ given by
$$x_i = \frac{\det A_i}{\det A}$$
where $A_i$ is the matrix obtained by replacing column $i$ of $A$ by $b$.
\end{theorem}

\begin{example}
Solve for $y$ only in the system of equations below
$$\begin{cases}
    x + y + z = 8 \\
    -x - 2y + 3z = 1 \\
    3x - 7y + 4z = 0
\end{cases}$$
\end{example}

\textit{Solution.}

Let $A$ be the coefficient matrix and $b$ be the column vector of the constants,
$$A = \bmat{
1 & 1 & 1 \\
-1 & -2 & 3 \\
3 & -7 & 4
}, \hspace{0.3cm} b = \bmat{
8 \\ 1 \\ 0
}, \hspace{0.3cm} A_2 = \bmat{1 & 8 & 1 \\ -1 & 1 & 3 \\ 3 & 0 & 4}
$$
By Cramer's Rule,
$$y = \frac{\det A_2}{\det A} = \frac{105}{39} = \boxed{\frac{35}{13}}$$